% ============================================================
%  Expanded Proofs
%  TU Berlin, Summer Semester 2025
% ============================================================

\documentclass[a4paper, 12pt]{article}

% ---- Font & layout (document-specific) ----------------------
\usepackage{mathpazo}
\linespread{1.05}
\usepackage{microtype}
\usepackage{geometry}
\geometry{margin=2.5cm, top=3cm, bottom=3cm}
\setlength{\parskip}{0.4em}
\setlength{\parindent}{1.5em}

% ---- Section headings ---------------------------------------
\usepackage{titlesec}
\titleformat{\section}
  {\large\bfseries}{\thesection.}{0.6em}{}
  [\vspace{0.3ex}\rule{\linewidth}{0.4pt}]
\titlespacing*{\section}{0pt}{2.5ex}{1.2ex}

% ---- Running headers ----------------------------------------
\usepackage{fancyhdr}
\pagestyle{fancy}
\fancyhf{}
\lhead{\small\itshape Expanded Proofs}
\rhead{\small\thepage}
\renewcommand{\headrulewidth}{0.4pt}
\fancypagestyle{plain}{%
  \fancyhf{}
  \rhead{\small\thepage}
  \renewcommand{\headrulewidth}{0pt}
}

% ---- Shared preamble ----------------------------------------
% ============================================================
%  Shared preamble — \input this from every notes document.
%  Do NOT put \documentclass, fonts, geometry, titlesec,
%  or fancyhdr here; those stay per-document.
% ============================================================

% ---- Encoding & typography ----------------------------------
\usepackage[utf8]{inputenc}
\usepackage[T1]{fontenc}
\usepackage{microtype}

% ---- Mathematics --------------------------------------------
\usepackage{amsmath, amsthm, amssymb}
\usepackage{mathtools}

% ---- Lists --------------------------------------------------
\usepackage{enumitem}


% ---- Colour & links -----------------------------------------
\usepackage{xcolor}
\usepackage{hyperref}
\hypersetup{colorlinks=true, linkcolor=black!60, citecolor=black!60, urlcolor=black!60}

% ---- Diagrams -----------------------------------------------
\usepackage{tikz-cd}

% ---- Theorem environments -----------------------------------
\theoremstyle{definition}
\newtheorem{example}{Example}
\newtheorem{remark}[example]{Remark}
\newtheorem{definition}[example]{Definition}

\theoremstyle{plain}
\newtheorem{theorem}[example]{Theorem}
\newtheorem{proposition}[example]{Proposition}
\newtheorem{lemma}[example]{Lemma}
\newtheorem{corollary}[example]{Corollary}
\newtheorem{claim}[example]{Claim}

% ---- Operators ----------------------------------------------
\DeclareMathOperator{\Spec}{Spec}
\DeclareMathOperator{\Proj}{Proj}
\DeclareMathOperator{\Gal}{Gal}
\DeclareMathOperator{\GL}{GL}
\DeclareMathOperator{\Lie}{Lie}
\DeclareMathOperator{\Tor}{Tor}
\DeclareMathOperator{\Char}{char}
\DeclareMathOperator{\height}{ht}
\DeclareMathOperator{\V}{\textbf{V}}

\newcommand{\Aut}{\mathrm{Aut}}
\newcommand{\Hom}{\mathrm{Hom}}
\newcommand{\Ext}{\mathrm{Ext}}
\newcommand{\EExt}{\mathcal{E}xt}
\newcommand{\Def}{\mathrm{Def}}
\newcommand{\Supp}{\mathrm{Supp}}

% ---- Blackboard bold ----------------------------------------
\renewcommand{\AA}{\mathbb{A}}
\newcommand{\GG}{\mathbb{G}}
\newcommand{\PP}{\mathbb{P}}
\newcommand{\QQ}{\mathbb{Q}}
\newcommand{\RR}{\mathbb{R}}
\newcommand{\CC}{\mathbb{C}}
\newcommand{\NN}{\mathbb{N}}
\newcommand{\ZZ}{\mathbb{Z}}

% ---- Calligraphic / fraktur ---------------------------------
\newcommand{\cO}{\mathcal{O}}   % structure sheaf
\newcommand{\cM}{\mathcal{M}}
\newcommand{\cB}{\mathcal{B}}
\newcommand{\cT}{\mathcal{T}}
\newcommand{\fX}{\mathfrak{X}}

% ---- Misc ---------------------------------------------------
\newcommand{\kk}{k}


% ---- Title --------------------------------------------------
\title{\textbf{Expanded Proofs}}
\author{Dominic Bunnett}
\date{Last updated: \today}

% =============================================================
\begin{document}
\maketitle
\tableofcontents
\newpage

% =============================================================
\section*{Introduction}
Expanded and detailed proofs of results from the Schemes Reading Course,
following Mumford's \emph{The Red Book of Varieties and Schemes}.

% =============================================================
\section{Chapter 2, Section 2}

\begin{theorem}[Adjunction isomorphism]
  Let $X$ be a scheme and $A$ a commutative ring. There is a natural isomorphism
  \[
    \Hom_{\mathrm{Sch}}(X,\, \Spec A)
    \;\cong\;
    \Hom_{\mathrm{Ring}}(A,\, \Gamma(X, \cO_X)).
  \]
\end{theorem}

\begin{proof}
  \begin{enumerate}
    \item \textbf{A ring map determines points.}
      Given a ring homomorphism $\phi \colon A \to \Gamma(X, \cO_X)$, we
      obtain a set map on underlying points: a prime $\mathfrak{p} \in X$
      determines a prime $\phi^{-1}(\mathfrak{p}) \in \Spec A$ via the
      composite $A \xrightarrow{\phi} \Gamma(X,\cO_X) \to \cO_{X,\mathfrak{p}}
      \to \kappa(\mathfrak{p})$, whose kernel is $\phi^{-1}(\mathfrak{p})$.

    \item \textbf{Reduce to affine opens.}
      Since $\Spec A$ is affine, a morphism of schemes $f \colon X \to \Spec A$
      is completely determined by its restrictions to an affine open cover
      $\{U_i\}$ of $X$ (together with the compatibility conditions on overlaps).
      It therefore suffices to construct the bijection on each affine open
      $U = \Spec B \subseteq X$, using the restricted ring map
      $A \xrightarrow{\phi} \Gamma(X,\cO_X) \xrightarrow{\mathrm{res}} B$.

    \item \textbf{Affine case.}
      For $X = \Spec B$ the required bijection is
      \[
        \Hom_{\mathrm{Sch}}(\Spec B,\, \Spec A)
        \;\cong\;
        \Hom_{\mathrm{Ring}}(A,\, B),
      \]
      which is the standard correspondence between ring maps $A \to B$ and
      scheme maps $\Spec B \to \Spec A$ \cite{Mumford}. Applying this to each
      affine open and observing that the resulting local morphisms glue
      (the compatibility on overlaps follows from naturality of the
      restriction maps in $\cO_X$) yields the global bijection. \qedhere
  \end{enumerate}
\end{proof}

% =============================================================
\bibliographystyle{alpha}
\begin{thebibliography}{9}

\bibitem{Mumford}
  D.~Mumford,
  \emph{The Red Book of Varieties and Schemes},
  Lecture Notes in Mathematics 1358,
  Springer, 2nd ed., 1999.

\end{thebibliography}

\end{document}
